\documentclass[exercice,a4]{thedocument}%
\usepackage[np]{numprint}%
\usepackage{enumitem}%
\startdatakeys
    \source{25GENMATPO3}
    \BO{2026}
    \cycle{4}
    \competencesDuSocle{CA,RA}
    \objectifsApprentissage{C4A1aa,C4A2ad,C4B2cb,C4D2aa,C4D3ac}
    \motsCles{fonction affine,image,antécédent,proportionnalité,graphique}
\enddatakeys
\begin{document}%
    \def\euro{\texteuro}%

\medskip

On veut poser du carrelage sur le sol intérieur d'une maison.

\medskip

Le carreleur A fait payer 80~\euro{} par m$^2$.

Le carreleur B fait payer 60~\euro{} par m$^2$ auxquels il faut ajouter 970~\euro{} pour la mise en place du chantier.

\medskip

\begin{enumerate}
\item Montrer que pour une surface dont l’aire est de 20 m$^2$, le prix est de \np{1600}~\euro{} avec le carreleur A et de \np{2170}~\euro{} avec le carreleur B.
\item Calculer le prix à payer pour une surface dont l’aire est 60 m$^2$ avec le carreleur A, puis avec le carreleur B.
\item On désigne par $x$ l’aire de la surface à carreler exprimée en m$^2$.

\begin{itemize}[label=$\bullet~$]
\item On appelle $f$ la fonction qui à l’aire à carreler en m$^2$ associe le prix en euros à payer avec
le carreleur A. On admet que $f$ est définie par $f(x) = 80x$.
\item On appelle $g$ la fonction qui à l’aire à carreler en m$^2$ associe le prix en euros à payer avec
le carreleur B. On admet que $g$ est définie par $g(x) = 60x + 970$.
\end{itemize}

	\begin{enumerate}
		\item Quelle est l'image de 70 par la fonction $f$ ?
		\item Quel est l'antécédent de \np{2400} par la fonction $f$ ?
		\item Sur le graphique fourni en ANNEXE, à rendre avec la copie, on a tracé la représentation graphique de la fonction $g$.

Tracer la représentation graphique de la fonction $f$ sur ce même graphique.
	\end{enumerate}
\item En utilisant le graphique fourni en ANNEXE, à rendre avec la copie, estimer l'aire maximale en m$^2$ que l'on peut carreler avec un budget de \np{2800}~\euro{} si l'on choisit le carreleur B.
\item Calculer l’aire en m$^2$ pour laquelle on paie exactement le même prix avec le carreleur A et le carreleur B.
\end{enumerate}

\bigskip
\showthesource
\end{document}
